% Research Paper Format for Chemistry Lab Reports
% Based on ACS formatting guidelines
% Use for: Labs with Abstract [+4+], Introduction [+6+], Experimental Details [+10+], Discussion [+18+], References [+2+]

\documentclass[
	10pt,
	twocolumn,
	letterpaper,
]{article}
\usepackage[
	top=0.75in,
	bottom=1in,
	left=0.75in,
	right=0.75in,
	columnsep=0.25in,
	headheight=14pt,
]{geometry}

\usepackage[utf8]{inputenc}
\usepackage[T1]{fontenc}
\usepackage{mathpazo}
\usepackage{microtype}
\usepackage{graphicx}
\usepackage{booktabs}
\usepackage{amsmath,amssymb}
\usepackage{float}

% Chemistry packages
\usepackage[version=4]{mhchem}
\usepackage{siunitx}
\sisetup{
	separate-uncertainty=true,
	range-phrase=--,
	range-units=single,
	per-mode=fraction,
	fraction-function=\frac
}
\usepackage{chemfig}

% Citations
\usepackage[
	backend=biber,
	style=chem-acs,
	sorting=none,
]{biblatex}
\addbibresource{references.bib}

% Formatting
\usepackage{xcolor}
\usepackage{hyperref}
\hypersetup{
	colorlinks=true,
	linkcolor=black,
	citecolor=blue!70!black,
	urlcolor=blue!70!black,
}
\usepackage{xurl}

\AtEveryCite{\color{blue!70!black}}

\usepackage{titlesec}
\titleformat{\section}{\large\bfseries\sffamily}{\thesection.}{0.5em}{}
\titleformat{\subsection}{\normalsize\bfseries\sffamily}{\thesubsection.}{0.5em}{}
\titleformat{\subsubsection}{\normalsize\itshape}{\thesubsubsection.}{0.5em}{}
\titlespacing*{\section}{0pt}{1.5ex plus 0.5ex minus 0.2ex}{1ex plus 0.2ex}
\titlespacing*{\subsection}{0pt}{1.2ex plus 0.3ex minus 0.1ex}{0.8ex plus 0.1ex}
\titlespacing*{\subsubsection}{0pt}{1ex plus 0.2ex minus 0.1ex}{0.5ex plus 0.1ex}

\usepackage{fancyhdr}
\pagestyle{fancy}
\fancyhf{}
\renewcommand{\headrulewidth}{0.4pt}
\fancyhead[L]{\small\textit{Experiment Title}}
\fancyhead[R]{\small\textit{Your Name}}
\fancyfoot[C]{\thepage}

\setlength{\parindent}{1em}
\setlength{\parskip}{0pt}

\usepackage{caption}
\captionsetup{
	font=small,
	labelfont=bf,
	justification=justified,
	singlelinecheck=false,
	skip=6pt,
}

\usepackage{balance}

% ============================================
% STUDENT INFORMATION - EDIT THESE VALUES
% ============================================
\newcommand{\reporttitle}{Experiment Title}
\newcommand{\reportauthor}{Your Name}
\newcommand{\reportpartner}{Lab Partner Name}
\newcommand{\reportinstructor}{Instructor Name}
\newcommand{\reportcourse}{Course Name}
\newcommand{\reportsection}{Section}
\newcommand{\reportinstitution}{Institution Name}
\newcommand{\reportdate}{\today}

\begin{document}

\twocolumn[
	\begin{@twocolumnfalse}
		\centering
		\vspace{0.5cm}
		{\Large\bfseries \reporttitle \par}
		\vspace{0.4cm}
		{\normalsize \reportauthor\textsuperscript{*} and \reportpartner \par}
		\vspace{0.2cm}
		{\small\itshape \reportinstitution, \reportcourse\ (Section \reportsection) \par}
		{\small\itshape Instructor: \reportinstructor \quad $\cdot$ \quad \reportdate \par}
		\vspace{0.6cm}

		\noindent\rule{\textwidth}{0.4pt}
		\vspace{0.3cm}
		\begin{minipage}{0.92\textwidth}
			\small
			\textbf{Abstract:}
			% ABSTRACT GOES HERE (≤200 words)
			% Summarize:
			% 1. Purpose of the experiment
			% 2. Methods used
			% 3. Key results
			% 4. Main conclusions
			% Keep it concise and factual
		\end{minipage}
		\vspace{0.3cm}
		\noindent\rule{\textwidth}{0.4pt}
		\vspace{0.5cm}
	\end{@twocolumnfalse}
]

\section{Introduction}

% INTRODUCTION GOES HERE
% Include:
% - Background information on the topic
% - Problem statement
% - Literature review (with citations \cite{key})
% - Objectives of this experiment
% - Relevant chemical principles

\vspace{1em}

% Example citation format: \cite{pubchem_benzoic_acid}

\section{Experimental Details}

\subsection{Materials and Apparatus}
\begin{itemize}
	\item Chemical amounts and purities
	\item Equipment used
	\item Instrumentation
\end{itemize}

\subsection{Procedure}

\subsubsection{Day 1: [Procedure Name]}

% Describe the procedure step-by-step
% Include masses, volumes, temperatures, times

\subsubsection{Day 2: [Procedure Name]}

% Describe analysis procedures

\section{Results}

\subsection{Experimental Data}

\begin{table}[htbp]
	\centering
	\caption{Descriptive title for data table.}
	\label{tab:data}
	\begin{tabular}{@{}lr@{}}
		\toprule
		\textbf{Measurement} & \textbf{Value} \\
		\midrule
		Measurement 1 & \SI{value}{unit} \\
		Measurement 2 & \SI{value}{unit} \\
		Measurement 3 & \SI{value}{unit} \\
		\bottomrule
	\end{tabular}
\end{table}

\subsection{Calculations}

\subsubsection{[Calculation Name]}

% Show all calculations with proper equations
% Use \begin{equation} for numbered equations
% Use \begin{align*} for multi-step calculations

Example:
\begin{equation}
	\%\ \text{Yield} = \frac{m_{\text{product}}}{m_{\text{theoretical}}} \times 100\%
	\label{eq:yield}
\end{equation}

\section{Discussion}

% DISCUSSION GOES HERE
% Include:
% - Interpretation of results
% - Comparison to literature values
% - Explanation of observations
% - Error analysis (sources of error, their effects, mitigation)
% - Relation to research questions

% Integrate answers to any questions from the handout
% Use citations: \cite{reference_key}

\subsection{[Subtopic 1]}

% Organize discussion by subtopics

\subsection{Error Analysis}

% Identify sources of error
% Explain how each error affected results
% Suggest improvements

\section{Conclusion}

% CONCLUSION GOES HERE
% Summarize:
% - Brief summary of the experiment
% - Main findings (quantitative results)
% - Success assessment
% - Skills learned

\section*{References}
\addcontentsline{toc}{section}{References}

\printbibliography[heading=none]

\balance

\end{document}
